% !TeX program = xelatex
% !TeX encoding = UTF-8
\documentclass[UTF8,a4paper,11pt]{ctexart}

\usepackage[a4paper,margin=2.3cm]{geometry}
\usepackage{amsmath,amssymb,bm,mathtools}
\usepackage{booktabs,tabularx,array,longtable}
\usepackage{graphicx}
\graphicspath{{./图形/}{./回归结果/ep/RLM_H1_独立同方差/}}
\usepackage{float}
\usepackage{xcolor}
\usepackage{enumitem}
\usepackage{caption}
\usepackage{hyperref}
\usepackage{url}

\newcolumntype{Y}{>{\raggedright\arraybackslash}X}
\newcolumntype{R}{>{\raggedleft\arraybackslash}X}
\setlength{\parindent}{2em}
\setlength{\parskip}{0.2em}
\setlength{\tabcolsep}{5pt}
\renewcommand{\arraystretch}{1.2}
\urlstyle{same}
\hypersetup{
  colorlinks=true,
  linkcolor=blue!55!black,
  urlcolor=blue!55!black,
  pdftitle={EP/PE 统计分析与因子测试计划}
}

\title{\textbf{EP/PE 统计分析与因子测试计划}}
\author{}
\date{2026 年 9 月 14 日}

\begin{document}
\maketitle

\begin{abstract}
本文汇总 EP/PE 因子的计算口径、描述性统计和申万一级行业 PE 分布，
并给出 EP、PB 等估值因子的测试计划。当前已经完成因子结果验收、描述性分析
和 EP 月度横截面 Huber RLM；IC、分层回测、OLS/WLS 对照、稳健性及交易成本
检验仍属于后续计划。
\end{abstract}

\tableofcontents
\clearpage

\section{第五周工作：EP、PB 等因子测试计划}

在完成数据完整性、准确性和时效性比较后，下一步将检验估值因子是否能够
稳定预测股票未来收益。测试依次包括测试前验收、截面回归、IC 检验、分层回测
和稳健性检验。

处理前再次明确因子生产口径：不根据 ST/PT 标记或次新股标记筛除股票，
次新标记只用于记录；主动置为缺失值的掩码仅来自停牌状态。此外，在行情表
原始股票代码上排除 50 只 \texttt{900xxx.BJ} 上海 B 股和 39 只
\texttt{200/201xxx.SZ} 深圳 B 股。此前两库数据比较未筛除这 39 只深圳 B 股，
将其计入了深市 A 股，本次计算已修正该口径。

因子生产口径与测试股票池口径应分开管理。本次因子横截面和后续测试池均按
行情代码口径保留沪深 A 股及北交所股票；不筛除 ST/PT 和上市未满一年的股票。
所有未来可能增加的筛选标记均须为历史时点可用数据，不能回写或重算已经交付
的原始因子。北交所股票在 2025 年发生六位代码切换；本次利用休市前后的收盘价、
前收盘价和总股本生成 242 对一一映射，日截面继续保留行情表当日代码，财报连接
和跨期收益则按统一证券身份处理。

\subsection{研究对象与因子方向}

主测试因子采用盈利收益率：
\begin{equation}
  \mathrm{EP}_{i,t}
  =\frac{\mathrm{NP}^{\mathrm{TTM}}_{i,t}}{\mathrm{MV}_{i,t}},
\end{equation}
其中，$\mathrm{NP}^{\mathrm{TTM}}_{i,t}$ 为交易日 $t$ 当时已经公开的 TTM
归母净利润，$\mathrm{MV}_{i,t}$ 为当日总市值。EP 越大表示估值越低，预期未来
收益可能越高。

交付的原始 PE 为 EP 的倒数，负利润对应负 PE：
\begin{equation}
  \mathrm{PE}^{\mathrm{raw}}_{i,t}
  =\frac{\mathrm{MV}_{i,t}}{\mathrm{NP}^{\mathrm{TTM}}_{i,t}}
  =\frac{1}{\mathrm{EP}_{i,t}}.
\end{equation}
传统市盈率测试仅在 $\mathrm{NP}^{\mathrm{TTM}}_{i,t}>0$ 时进行。由于 PE 在净利润
接近零时容易产生极端值，并且负 PE 不能被解释为正常估值倍数，本文以 EP
作为主因子，以正利润样本的 PE 作为稳健性检验。

PB 因子可沿用相同测试框架。为使因子方向与 EP 一致，主测试可使用
$\mathrm{BP}=1/\mathrm{PB}$，此时 BP 越高代表估值越低、预期收益越高；若直接
使用 PB，则预期收益方向相反。

\subsection{测试前验收}

虽然数据清洗和因子计算已经完成，但在正式测试前仍需对最终因子表进行一次
验收，防止数据问题传递至回测结果。

\begin{enumerate}[leftmargin=2.7em]
  \item 在 MAD 去极值和标准化之前确定沪深纯 A 股和北交所测试池，排除
        89 只沪深 B 股；\textbf{已实现}。
  \item 保留退市股票在实际退市前的历史样本，避免幸存者偏差；\textbf{已实现}。
  \item 财务数据的公告时间不得晚于因子形成时间；同一报告期按公告时间确定
        版本，新公告生效后不允许旧公告覆盖；\textbf{已实现}。
  \item 保留总市值、TTM 归母净利润、报告期、公告时间、原始 EP/PE 和
        标准化因子值，确保异常结果可以追溯；\textbf{已实现}。
  \item 检查每日因子覆盖率，并从覆盖率稳定后的首个月末开始正式测试；
        \textbf{已实现}。EP 覆盖率连续 3 个自然月末不低于 95\%，首次确认日为
        2020 年 7 月 31 日。
\end{enumerate}

当前行情表没有历史 ST 标记，并且缺少完整的官方上市日期，因此当前测试不能
可靠执行这两类历史筛选，也不能使用当前股票名称或未来状态回填历史。

\subsection{样本与前瞻收益率}

测试频率采用月度：每个月最后一个交易日形成因子值，下个月第一个交易日按
复权收盘价建仓，再下一个月第一个交易日按复权收盘价退出。前瞻收益率定义为
\begin{equation}
  r^{\mathrm{fwd}}_{i,t+1}
  =\frac{P^{\mathrm{adj}}_{i,\mathrm{exit}}}
  {P^{\mathrm{adj}}_{i,\mathrm{entry}}}-1.
\end{equation}
复权价格用于消除分红、拆股和合股等公司行为造成的价格机械变化，使计算结果
反映投资者在两个固定交易端点之间的实际持有收益。

\subsection{分期截面回归}

每个月对股票下一期收益率进行一次截面回归：
\begin{equation}
  r^{\mathrm{fwd}}_{i,t+1}
  =\alpha_t+\beta_t X_{i,t}
  +\gamma_t\ln(\mathrm{MV}_{i,t})
  +\sum_{k=1}^{K-1}\delta_{k,t}I_{i,k,t}
  +\varepsilon_{i,t+1}.
\end{equation}
其中，$X_{i,t}$ 为标准化 EP，$I_{i,k,t}$ 为股票所属行业的哑变量，
$\ln(\mathrm{MV}_{i,t})$ 用于控制市值效应，$\beta_t$ 为当期因子收益。

主结果采用 RLM 稳健回归，以降低异常收益率对结果的影响；后续再使用 OLS
和以市值平方根加权的 WLS 进行对照。统计指标包括：
\begin{itemize}[leftmargin=2.5em]
  \item 因子收益均值及累计因子收益；
  \item 单月 RLM 系数的 H1 标准误及大样本 $z$ 值；
  \item 月度因子收益均值的 IID 标准误、$t$ 值和 $p$ 值，本阶段不使用
        Newey--West HAC；
  \item $\beta_t>0$ 的月份比例；
  \item 单期回归 $|z|\geq2$ 的月份比例；
  \item RLM 与后续 OLS/WLS 结果的一致性。
\end{itemize}

对于 EP 或 BP 因子，预期回归系数为正；若直接使用 PE 或 PB，则预期系数
为负。EP 的新口径 RLM 已完成，结果见第三部分；OLS/WLS 对照尚未运行。

\subsection{Rank IC 检验}

首先将 EP 对行业和对数市值进行截面回归，使用回归残差
$X^{\perp}_{i,t}$ 作为行业和市值中性化后的因子暴露。Rank IC 定义为
\begin{equation}
  IC_t=\operatorname{Spearman}
  \left(X^{\perp}_{i,t},r^{\mathrm{fwd}}_{i,t+1}\right),
  \qquad
  IR=\frac{\overline{IC}}{\sigma(IC)}.
\end{equation}
IC 检验应报告 IC 均值、标准差、$IC_t>0$ 的月份比例、$|IC_t|>0.02$ 的月份
比例以及 IR。还应逐年统计上述指标，检查因子的预测能力是否集中在少数年份。
本阶段不执行 Newey--West HAC 推断。

\subsection{行业内分层回测}

每个调仓日在各一级行业内部按照 EP 从低到高划分为五组，再将各行业相同序号
的组别合并为全市场组合。主结果采用市值加权，同时报告等权结果。多空组合收益
定义为
\begin{equation}
  R^{\mathrm{LS}}_t=R^{Q_5}_t-R^{Q_1}_t,
\end{equation}
其中，$Q_5$ 为 EP 最高、估值最低的一组，$Q_1$ 为 EP 最低的一组。

各组及多空组合需要统计年化收益、累计收益、基准超额收益、年化波动率、
夏普比率、最大回撤、信息比率、胜率和换手率。通过五组收益排列以及组序号
与收益之间的 Spearman 相关系数检验因子单调性。基准优先选择中证全指；
若缺少指数数据，可使用测试股票池的市值加权收益作为内部基准。

\subsection{交易成本与稳健性检验}

根据组合换手率计算扣除交易成本后的收益，并设置不同单边成本情景。此外，
还需要完成以下检验：
\begin{enumerate}[leftmargin=2.7em]
  \item 分年度以及市场上涨、下跌阶段统计回归、IC 和多空收益；
  \item 分别在大、中、小市值股票中进行测试；
  \item 比较全样本 EP 与仅包含正利润公司的传统 PE 样本；
  \item 比较等权与市值加权、五组与十组、T+1 与 T+2 生效结果；
  \item 比较不同 MAD 参数下测试结论是否稳定；
  \item 固定行情和股票池，分别使用 Wind 与 base 财报构造因子，比较两套
        因子的覆盖率、秩相关性、IC 和多空收益；
  \item 将最后一段时间保留为样本外区间，不根据样本外结果调整参数。
\end{enumerate}

\subsection{综合评价标准}

\begin{table}[H]
  \centering
  \caption{PE/EP 因子综合评价标准}
  \begin{tabularx}{\textwidth}{lYY}
    \toprule
    评价维度 & 核心指标 & 判断方法 \\
    \midrule
    有效性 & 因子收益、单月 H1 $z$ 值及全期 IID $t$ 值
      & 收益方向符合预期，且具有统计与经济意义 \\
    预测能力 & Rank IC、IR
      & IC 方向稳定，均值不依赖少数极端月份 \\
    单调性 & 五组收益及 $Q_5-Q_1$
      & 各组收益大体单调，多空收益方向正确 \\
    稳定性 & 年度及市场状态分组
      & 结果不依赖少数月份或单一年份 \\
    可交易性 & 换手率及扣费后收益
      & 扣除合理交易成本后仍具有正收益 \\
    数据稳健性 & Wind 与 base 测试结果
      & 核心结论不因财报数据来源而反转 \\
    \bottomrule
  \end{tabularx}
\end{table}

只有当截面回归、Rank IC 和分层回测的方向基本一致，样本外结果稳定，并且
扣除交易成本后仍有收益时，才能认为 PE/EP 因子具有较可靠的实际选股能力。
所有阈值和参数应在查看样本外结果之前固定，避免结果导向调参。

\subsection{实施顺序与输出结果}

\begin{equation*}
\begin{aligned}
  &\text{测试前验收}
  \longrightarrow \text{前瞻收益构造}
  \longrightarrow \text{RLM/OLS/WLS 截面回归} \\
  &\hspace{6em}\longrightarrow \text{Rank IC}
  \longrightarrow \text{行业内分层}
  \longrightarrow \text{稳健性与成本检验}.
\end{aligned}
\end{equation*}

最终形成因子数据验收表、截面回归汇总表、IC 汇总表、五组及多空组合绩效表、
年度稳定性表和两库因子测试对比表，作为判断该因子是否能够进入后续多因子模型
的依据。当前状态为：数据验收、描述性统计和 EP 月度横截面 RLM 已完成；其余
测试待执行。

\clearpage
\section{EP/PE 因子结果统计}

\subsection{样本与计算口径}

样本期为 2020 年 1 月 2 日至 2025 年 12 月 31 日，共 1,455 个交易日。
在行情表原始代码上排除 50 只 \texttt{900xxx.BJ} 上海 B 股和 39 只
\texttt{200/201xxx.SZ} 深圳 B 股后，样本包含 7,081,800 条股票日记录和
5,927 个全样本期唯一行情代码标识。其中沪市 2,382 个、深市 3,006 个、
北交所 539 个。这里的 539 是全样本期代码标识数，不是北交所公司数：其中包括
251 个 \texttt{43/83/87xxxx.BJ} 历史代码和 288 个 \texttt{920xxx.BJ} 新代码；
2025 年代码切换使同一证券的新旧代码被重复计数。2025 年 12 月 31 日行情截面
共有 288 只北交所股票，均使用 \texttt{920xxx.BJ} 代码。实际排除的证券只有
上述 89 只沪深 B 股，普通北交所股票全部保留。ST/PT 股票不作筛选，次新股标记
仅用于记录，只有停牌日被主动置为缺失值。

原始因子定义为
\begin{equation}
  \mathrm{EP}_{i,t}
  =\frac{\text{TTM 归母净利润}_{i,t}}{\text{总市值}_{i,t}},
  \qquad
  \mathrm{PE}_{i,t}=\frac{1}{\mathrm{EP}_{i,t}}.
\end{equation}
行情日 $t$ 仅使用截至当日结束前已经可获得的财务信息，并将因子用于下一交易日。
原始 EP 和 PE 各有 6,732,630 条有效观测，有效率均为 95.07\%，点时未来数据
违规为 0。北交所代码切换后，行情代码仍保存为 \texttt{920xxx.BJ}，但通过
\texttt{financial\_code6} 连接同一证券旧代码下的财报历史。

\subsection{数据验收结果}

\begin{table}[H]
  \centering
  \caption{因子数据验收结果}
  \begin{tabularx}{0.92\textwidth}{Y R}
    \toprule
    验收项目 & 结果 \\
    \midrule
    排除 B 股 & 89 只：上海 B 股 50 只、深圳 B 股 39 只 \\
    北交所新旧代码映射 & 242 对，价格连续且一一对应 \\
    忽略旧公告回退记录 & 3,881 条 \\
    有效财报版本 & 262,984 条 \\
    逐行因子追溯记录 & 7,081,800 条 \\
    点时未来数据违规 & 0 条 \\
    覆盖率稳定规则 & 非停牌样本连续 3 个自然月末不低于 95\% \\
    正式测试起点 & 2020 年 7 月 31 日 \\
    \bottomrule
  \end{tabularx}
\end{table}

逐行追溯表保存行情代码、统一证券身份、财务连接代码、市场、市值、当日交易状态、
停牌标记、上市日期可用性、实际采用的报告期、公告时间、版本时间、TTM 归母净利润、
原始 EP/PE、MAD 处理值和 Z 标准化值。历史 ST 和官方上市日期没有源表，相关可用性
及“不筛选”状态分别记录在追溯表和 \texttt{point\_in\_time\_field\_audit.csv} 中。

\subsection{描述性统计}

\begin{table}[H]
  \centering
  \caption{原始 EP/PE 描述性统计}
  \label{tab:ep-pe-descriptive}
  \begin{tabular}{lrr}
    \toprule
    统计量 & EP & PE \\
    \midrule
    有效观测数 & 6,732,630 & 6,732,630 \\
    均值 & 0.006925 & $-13.3974$ \\
    标准差 & 0.174645 & 12,684.6366 \\
    1\% 分位数 & $-0.416115$ & $-592.1720$ \\
    中位数 & 0.021148 & 26.2719 \\
    99\% 分位数 & 0.174537 & 751.4696 \\
    负值占比 & 21.36\% & 21.36\% \\
    MAD 截断占比 & 11.07\% & 14.16\% \\
    \bottomrule
  \end{tabular}
\end{table}

最新截面为 2025 年 12 月 31 日，共有 5,438 条有效观测；EP 中位数为 0.01315，
PE 中位数为 28.18 倍，负 EP/PE 占比为 28.26\%。经过逐日 MAD 去极值和 Z
标准化后，EP 和 PE 各保留 6,732,629 条有效结果。

全样本中，PE 的最小值和最大值分别约为 $-3.73\times10^6$ 倍和
$5.39\times10^5$ 倍，说明 EP 接近零时倒数变换会形成明显长尾。因此，PE
更适合使用中位数、分位数或去极值结果进行解释，EP 则更适合直接进行横截面比较。

\subsection{分行业原始 PE}

申万一级行业分析使用 2025 年 12 月 31 日最新截面，按行业汇总未经 MAD 去极值
和 Z 标准化的原始 PE。为避免接近零的 EP 使 PE 均值被极端值主导，柱高采用
行业内原始 PE 的中位数。亏损公司的负 PE 仍保留在统计中。

行业中位 PE 最高的三个行业分别为通信（57.55 倍）、综合（53.22 倍）和国防军工
（52.75 倍）。房地产的行业中位 PE 为 $-1.89$ 倍，其负 PE 比例为 62.00\%；
负 PE 表示行业内亏损观测较多，不应解释为负估值倍数。图中红色虚线表示全市场
原始 PE 中位数 28.18 倍。

\begin{figure}[H]
  \centering
  \includegraphics[width=0.96\textwidth]{industry_pe_latest.png}
  \caption{2025 年 12 月 31 日申万一级行业原始 PE 中位数。样本包含 5,438 条
  PE 和行业分类均有效的股票记录；红色虚线为全市场原始 PE 中位数 28.18 倍，
  红色柱表示行业中位 PE 为负。}
  \label{fig:industry-pe}
\end{figure}

\clearpage
\section{EP 横截面 RLM 回归}

\subsection{横截面线性模型}

在每个因子形成日 $t$，设横截面内共有 $N$ 只股票和 $K$ 个行业。下一持有期
收益写成 $N\times1$ 列向量：
\begin{equation}
  \bm{Y}_t=
  \begin{bmatrix}
    r_{1,t+1} & r_{2,t+1} & \cdots & r_{N,t+1}
  \end{bmatrix}^{\mathsf T}.
\end{equation}

EP 因子暴露和市值暴露分别写成
\begin{equation}
  \bm{x}_t=
  \begin{bmatrix}x_{1,t}&x_{2,t}&\cdots&x_{N,t}\end{bmatrix}^{\mathsf T},
  \qquad
  \bm{m}_t=
  \begin{bmatrix}m_{1,t}&m_{2,t}&\cdots&m_{N,t}\end{bmatrix}^{\mathsf T}.
\end{equation}
其中，$x_{i,t}$ 是股票 $i$ 在形成日的标准化 EP，$m_{i,t}$ 是标准化后的对数
总市值。行业归属可写成 $N\times K$ 的哑变量矩阵。例如：
\begin{equation*}
\bm{I}_t=
\begin{array}{c|ccc}
  & \text{行业1} & \text{行业2} & \text{行业3} \\ \hline
  \text{股票1} & 1&0&0\\
  \text{股票2} & 0&1&0\\
  \text{股票3} & 0&0&1\\
  \text{股票4} & 1&0&0
\end{array}.
\end{equation*}

每只股票只属于一个一级行业，因此每一行只有一个元素为 1。为了同时估计截距
并避免行业哑变量陷阱，回归只放入 $K-1$ 个行业哑变量。设计矩阵和待估参数为
\begin{align}
  \bm{D}_t
  &=\begin{bmatrix}
      \bm{1}_N&\bm{x}_t&\bm{m}_t&\bm{I}_{1,t}&\cdots&\bm{I}_{K-1,t}
    \end{bmatrix},\\
  \bm{\theta}_t
  &=\begin{bmatrix}
      \alpha_t&\beta_{\mathrm{EP},t}&\beta_{\mathrm{Size},t}&
      \gamma_{1,t}&\cdots&\gamma_{K-1,t}
    \end{bmatrix}^{\mathsf T}.
\end{align}
因此当月横截面模型为
\begin{align}
  \bm{Y}_t&=\bm{D}_t\bm{\theta}_t+\bm{\varepsilon}_t,\\
  r_{i,t+1}
  &=\alpha_t+\beta_{\mathrm{EP},t}x_{i,t}
    +\beta_{\mathrm{Size},t}m_{i,t}
    +\sum_{k=1}^{K-1}\gamma_{k,t}I_{i,k,t}
    +\varepsilon_{i,t+1}.
\end{align}
这里的系数在同一个月内对所有股票共同估计，所以应写成 $\alpha_t$、
$\beta_{\mathrm{EP},t}$ 和 $\gamma_{k,t}$，而不是带股票下标的系数。
$\beta_{\mathrm{EP},t}$ 就是该月 EP 暴露对应的横截面因子收益。

\subsection{OLS 的估计与统计推断}

OLS 选择使残差平方和最小的参数：
\begin{equation}
  S(\bm{\theta}_t)
  =\sum_{i=1}^{N}(y_i-\bm{d}_i^{\mathsf T}\bm{\theta}_t)^2
  =(\bm{Y}_t-\bm{D}_t\bm{\theta}_t)^{\mathsf T}
   (\bm{Y}_t-\bm{D}_t\bm{\theta}_t).
\end{equation}
对 $\bm{\theta}_t$ 求一阶导数并令其为零：
\begin{equation}
  \frac{\partial S}{\partial\bm{\theta}_t}
  =-2\bm{D}_t^{\mathsf T}
  (\bm{Y}_t-\bm{D}_t\bm{\theta}_t)=0.
\end{equation}
得到正规方程及其形式解：
\begin{align}
  \bm{D}_t^{\mathsf T}\bm{D}_t\widehat{\bm{\theta}}_t
    &=\bm{D}_t^{\mathsf T}\bm{Y}_t,\\
  \widehat{\bm{\theta}}_t
    &=(\bm{D}_t^{\mathsf T}\bm{D}_t)^{-1}
      \bm{D}_t^{\mathsf T}\bm{Y}_t.
\end{align}

如果同时加入截距和全部 $K$ 个行业哑变量，则行业哑变量之和恒等于截距列，
$\bm{D}_t$ 不满列秩，$\bm{D}_t^{\mathsf T}\bm{D}_t$ 没有逆。因此本次使用
“截距加 $K-1$ 个行业哑变量”，并在每个月拟合前检查设计矩阵的秩。实际数值
计算不显式求逆，而使用稳定的线性代数求解器。

设参数总数为 $p=K+2$，OLS 残差和共同残差方差估计为
\begin{align}
  \widehat{\bm{\varepsilon}}_t
    &=\bm{Y}_t-\bm{D}_t\widehat{\bm{\theta}}_t,\\
  \widehat{\sigma}_t^2
    &=\frac{\widehat{\bm{\varepsilon}}_t^{\mathsf T}
      \widehat{\bm{\varepsilon}}_t}{N-p}.
\end{align}
在残差独立同方差的经典假设下：
\begin{equation}
  \widehat{\operatorname{Var}}_{\mathrm{OLS}}
  (\widehat{\bm{\theta}}_t)
  =\widehat{\sigma}_t^2
   (\bm{D}_t^{\mathsf T}\bm{D}_t)^{-1}.
\end{equation}
第 $j$ 个系数的标准误和统计量为
\begin{equation}
  SE(\widehat{\theta}_{j,t})
  =\sqrt{\left[
  \widehat{\operatorname{Var}}(\widehat{\bm{\theta}}_t)
  \right]_{jj}},
  \qquad
  t_{j,t}=\frac{\widehat{\theta}_{j,t}}{SE(\widehat{\theta}_{j,t})}.
\end{equation}

标准误衡量重复抽样时系数估计值的波动程度；统计量衡量当前系数距离零有多少个
标准误。经典 OLS 在正态、独立同方差等条件下使用自由度为 $N-p$ 的 Student
$t$ 分布：
\begin{equation}
  p_j=2\Pr\!\left(T_{N-p}>|t_{j,t}|\right),
  \qquad
  CI_{95\%}=\widehat{\theta}_{j,t}
  \pm t_{0.975,N-p}SE(\widehat{\theta}_{j,t}).
\end{equation}

\begin{table}[H]
  \centering
  \caption{大样本双侧检验常用临界值}
  \begin{tabular}{cc}
    \toprule
    显著性水平 & 临界值（近似） \\
    \midrule
    10\% & $|t|>1.65$ \\
    5\%  & $|t|>1.96$ \\
    1\%  & $|t|>2.58$ \\
    \bottomrule
  \end{tabular}
\end{table}

\subsection{为什么主回归使用 RLM}

股票收益横截面通常呈现尖峰厚尾，极端收益出现概率远高于正态分布。OLS 使用
残差平方和作为目标函数，会让极端残差的影响按平方增加。RLM 不直接删除这些
观测，而是对大残差连续降权，使其仍保留信息但不再主导系数。

RLM 的目标函数为
\begin{equation}
  \min_{\bm{\theta}_t}
  \sum_{i=1}^{N}
  \rho\!\left(
  \frac{y_i-\bm{d}_i^{\mathsf T}\bm{\theta}_t}{s_t}
  \right),
\end{equation}
其中 $s_t$ 是使用 MAD 估计的当月稳健残差尺度。本次采用 Huber 损失，
$c=1.345$：
\begin{equation}
\rho(u)=
\begin{cases}
  \frac{1}{2}u^2, & |u|\leq c,\\
  c|u|-\frac{1}{2}c^2, & |u|>c.
\end{cases}
\end{equation}
$c=1.345$ 使 Huber 估计在正态误差下保留约 95\% 的渐近效率，同时对厚尾异常值
具有较强抵抗力。对应的得分函数和 IRLS 权重为
\begin{align}
  \psi(u)=\rho'(u)
  &=\begin{cases}
      u,&|u|\leq c,\\
      c\operatorname{sign}(u),&|u|>c,
    \end{cases}\\
  w(u)=\frac{\psi(u)}{u}
  &=\begin{cases}
      1,&|u|\leq c,\\
      c/|u|,&|u|>c.
    \end{cases}
\end{align}

\subsubsection{IRLS 求解过程}

RLM 使用迭代加权最小二乘求解：
\begin{enumerate}[leftmargin=2.7em]
  \item 使用 OLS 得到初始系数：
  \begin{equation}
    \bm{\theta}^{(0)}
    =(\bm{D}^{\mathsf T}\bm{D})^{-1}\bm{D}^{\mathsf T}\bm{Y}.
  \end{equation}

  \item 根据初始系数计算残差和 MAD 稳健尺度：
  \begin{equation}
    \bm{e}^{(q)}=\bm{Y}-\bm{D}\bm{\theta}^{(q)},
    \qquad
    s^{(q)}=\operatorname{MAD}(\bm{e}^{(q)}).
  \end{equation}
  MAD 中使用的正态一致性常数约为 $0.6745$，使其在正态分布下与标准差处于
  相同尺度。

  \item 计算标准化残差和 Huber 权重：
  \begin{equation}
    u_i^{(q)}=\frac{e_i^{(q)}}{s^{(q)}},
    \qquad
    \bm{W}^{(q)}=\operatorname{diag}
    \left(w_1^{(q)},\ldots,w_N^{(q)}\right).
  \end{equation}

  \item 使用本轮权重进行一次加权最小二乘：
  \begin{equation}
    \bm{\theta}^{(q+1)}
    =(\bm{D}^{\mathsf T}\bm{W}^{(q)}\bm{D})^{-1}
      \bm{D}^{\mathsf T}\bm{W}^{(q)}\bm{Y}.
  \end{equation}

  \item 重新计算残差、MAD 尺度和权重，直至满足
  \begin{equation}
    \max_j|\theta_j^{(q+1)}-\theta_j^{(q)}|\leq10^{-8},
  \end{equation}
  或达到 100 次迭代上限。本次 64 个月度模型均在达到上限之前收敛。
\end{enumerate}

\subsubsection{H1 协方差与单月标准误}

为简化符号，以下省略月份下标。令
\begin{equation}
  u_i=\frac{y_i-\bm{d}_i^{\mathsf T}\bm{\theta}}{s}.
\end{equation}
RLM 的一阶条件为
\begin{equation}
  \bm{g}(\bm{\theta})
  =\sum_{i=1}^{N}\bm{d}_i\psi(u_i)=\bm{0}.
\end{equation}
在真实参数 $\bm{\theta}_0$ 附近进行一阶泰勒展开：
\begin{equation}
  \widehat{\bm{\theta}}-\bm{\theta}_0
  \approx
  s\bm{A}^{-1}
  \sum_{i=1}^{N}\bm{d}_i\psi(u_i),
  \qquad
  \bm{A}=\sum_{i=1}^{N}\psi'(u_i)\bm{d}_i\bm{d}_i^{\mathsf T}.
\end{equation}
利用 $\operatorname{Cov}(\bm{C}\bm{Z})
=\bm{C}\operatorname{Cov}(\bm{Z})\bm{C}^{\mathsf T}$，并在不同股票得分相互独立
的条件下，得到一般三明治协方差：
\begin{align}
  \widehat{\operatorname{Var}}(\widehat{\bm{\theta}})
  &\approx s^2\bm{A}^{-1}\bm{B}\bm{A}^{-1},\\
  \bm{B}
  &=\sum_{i=1}^{N}\psi(u_i)^2\bm{d}_i\bm{d}_i^{\mathsf T}.
\end{align}

H1 将残差得分部分与设计矩阵部分分开。定义
\begin{align}
  m&=\frac{1}{N}\sum_{i=1}^{N}\psi'(u_i),\\
  q&=\frac{1}{N-p}\sum_{i=1}^{N}\psi(u_i)^2,\\
  v&=\operatorname{Var}\!\left(\psi'(u_i)\right),\\
  k&=1+\frac{p}{N}\frac{v}{m^2}.
\end{align}
则 \texttt{statsmodels} 使用的 H1 协方差为
\begin{equation}
  \boxed{
  \widehat{\operatorname{Var}}_{H1}(\widehat{\bm{\theta}})
  =k^2\frac{s^2q}{m^2}
   (\bm{D}^{\mathsf T}\bm{D})^{-1}}
  \label{eq:h1-cov}
\end{equation}
EP 系数的单月 H1 标准误和统计量为
\begin{equation}
  SE_{H1}(\widehat\beta_{\mathrm{EP}})
  =\sqrt{\left[
  \widehat{\operatorname{Var}}_{H1}(\widehat{\bm{\theta}})
  \right]_{\mathrm{EP},\mathrm{EP}}},
  \qquad
  z_{H1}=\frac{\widehat\beta_{\mathrm{EP}}}
  {SE_{H1}(\widehat\beta_{\mathrm{EP}})}.
\end{equation}

当不截断任何残差时，$\psi(u)=u$、$\psi'(u)=1$、$m=1$ 且 $k=1$，
式~\eqref{eq:h1-cov} 退化为经典 OLS 协方差。Huber 截断后，大残差在
$\sum\psi(u_i)^2$ 中的贡献被限制；同时，$m<1$ 会补偿降权造成的信息损失。

H1 标准误对应的是大样本近似。对 M 估计量，在当月股票数量 $N$ 增大时：
\begin{equation}
  \sqrt{N}
  (\widehat{\bm{\theta}}-\bm{\theta}_0)
  \xrightarrow{d}\mathcal{N}(\bm{0},\bm{V}).
\end{equation}
因此在原假设 $H_0:\beta_{\mathrm{EP}}=0$ 下，$z_{H1}$ 近似服从标准正态分布。
这里的大样本是每个月约 4,000 至 5,400 只股票，而不是 64 个月。由于 Huber
权重和 MAD 尺度均由数据决定，RLM 不具备经典 OLS 有限样本 Student $t$ 分布
所要求的精确结构，因此单月统计量应称为大样本 $z$ 值。

“当月横截面参数的不确定性”是指：在当月因子、市值和行业暴露给定的情况下，
如果重新产生一组股票特质残差并再次拟合，$\widehat\beta_{\mathrm{EP},t}$ 会如何
波动。H1 标准误估计这种重复抽样波动，不回答 EP 系数在不同月份之间是否稳定。
H1 对异常残差稳健，但没有处理股票之间的横截面相关性。

\subsubsection{月度系数的全期汇总}

64 个月分别产生 64 个 $\widehat\beta_{\mathrm{EP},t}$、64 个 H1 标准误和
64 个 H1 $z$ 值。判断平均因子收益时，统计单位变为月份。本次暂时假设月度系数
相互独立且同方差：
\begin{align}
  \overline\beta_{\mathrm{EP}}
  &=\frac{1}{M}\sum_{t=1}^{M}\widehat\beta_{\mathrm{EP},t},\\
  SE_{\mathrm{IID}}(\overline\beta_{\mathrm{EP}})
  &=\frac{s_{\beta}}{\sqrt{M}},\\
  t_{\mathrm{IID}}
  &=\frac{\overline\beta_{\mathrm{EP}}}
  {SE_{\mathrm{IID}}(\overline\beta_{\mathrm{EP}})},\\
  p_{\mathrm{IID}}
  &=2\Pr\!\left(T_{M-1}>|t_{\mathrm{IID}}|\right).
\end{align}
这里的 $SE_{\mathrm{IID}}$ 是“EP 月均系数的时间序列 IID 标准误”，不是单月
H1 标准误。本次没有执行 Newey--West HAC。

\subsubsection{一日横截面收益率分布示例}

为直观展示 RLM 所处理的异常收益问题，选取 2025 年 12 月 29 日，对当日股票
计算未经截尾或去极值的复权收益率：
\begin{equation}
  r^{\mathrm{daily}}_{i,t}
  =\frac{P^{\mathrm{adj}}_{i,t}}
  {P^{\mathrm{adj}}_{i,t-1}}-1.
\end{equation}
该日原始行情有 5,545 条记录，按研究口径排除 79 条沪深 B 股记录和 11 条停牌
记录后，有效横截面包含 5,455 只股票。选择该日是为了使用接近样本末端、且没有
被上市首日数倍涨幅主导的完整交易日，日期选择不会改变正式回归样本。

\begin{table}[H]
  \centering
  \caption{2025 年 12 月 29 日横截面收益率分布}
  \begin{tabular}{lrr}
    \toprule
    统计量 & 实际分布 & 同均值、同标准差正态分布 \\
    \midrule
    均值 & $-0.2172\%$ & $-0.2172\%$ \\
    中位数 & $-0.4878\%$ & $-0.2172\%$ \\
    标准差 & 2.8354\% & 2.8354\% \\
    偏度 & 1.1009 & 0 \\
    超额峰度 & 12.8128 & 0 \\
    均值正负 0.5 个标准差以内 & 61.72\% & 38.29\% \\
    均值正负 3 个标准差以外 & 2.60\%（142 只） & 0.27\%（约 15 只） \\
    \bottomrule
  \end{tabular}
\end{table}

\begin{figure}[H]
  \centering
  \includegraphics[width=0.94\textwidth]{daily_return_distribution_20251229.png}
  \caption{2025 年 12 月 29 日股票横截面单日收益率分布。}
  \label{fig:daily-return-distribution}
\end{figure}

均值、方差、偏度和峰度分别描述分布的一阶至四阶矩。偏度描述分布相对于均值
的不对称程度；超额峰度反映相对于正态分布的尾部厚度。该日超额峰度明显大于零，
且三倍标准差之外的实际观测远多于正态分布预期，说明采用稳健回归具有实际意义。

\subsection{固定测试口径}

\begin{itemize}[leftmargin=2.5em]
  \item 因子形成日：每个月最后一个交易日；
  \item 建仓日：下个月第一个交易日，以 \texttt{close\_adj} 复权收盘价建仓；
  \item 退出日：再下一个月第一个交易日，以 \texttt{close\_adj} 复权收盘价退出；
  \item 因子：形成日经过 MAD 去极值和 Z 标准化的 EP；
  \item 控制变量：形成日横截面标准化后的对数总市值，以及 $K-1$ 个申万一级
        行业哑变量；
  \item 样本不筛除 ST/PT 或次新股，固定成交端点停牌或价格缺失时前瞻收益
        记为缺失。
\end{itemize}

测试模型为
\begin{equation}
  r^{\mathrm{fwd}}_{i,m+1}
  =\alpha_m+\beta_{\mathrm{EP},m}\mathrm{EPZ}_{i,m}
   +\beta_{\mathrm{Size},m}\mathrm{SizeZ}_{i,m}
   +\sum_{k=1}^{K-1}\gamma_{k,m}I_{i,k,m}
   +\varepsilon_{i,m+1}.
\end{equation}
RLM 使用 HuberT（$c=1.345$）、MAD 残差尺度和 H1 协方差，单月系数使用 H1
标准误计算大样本 $z$ 值。月度系数序列的全期和年度汇总暂按独立同方差假设
计算普通均值 $t$ 检验，不执行 Newey--West HAC。

\subsection{样本与拟合情况}

共得到 64 个完整持有期，因子形成日从 2020 年 7 月 31 日至 2025 年 10 月 31 日，
首次建仓日为 2020 年 8 月 3 日，最后退出日为 2025 年 12 月 1 日。控制变量和
行业均有效的观测为 313,200 条，其中 312,177 条具有有效固定端点收益。单月最终
样本数介于 3,812 和 5,400 之间，中位数为 5,011。全部 64 个月度模型均收敛。

\subsection{全期结果}

\begin{table}[H]
  \centering
  \caption{EP 月度横截面 RLM 全期汇总}
  \begin{tabular}{lr}
    \toprule
    指标 & 结果 \\
    \midrule
    回归月份 & 64 \\
    EP 月均系数 & 0.1623\% \\
    EP 系数中位数 & 0.1465\% \\
    EP 系数标准差 & 0.8964\% \\
    EP 月均系数的时间序列 IID 标准误 & 0.1120\% \\
    EP 月均系数的 IID $t$ 值 & 1.4487 \\
    EP 月均系数的 IID $p$ 值 & 0.1524 \\
    EP 系数为正的月份 & 53.12\% \\
    单月 $|z_{H1}|\geq2$ 的比例 & 82.81\% \\
    正向且 $z_{H1}\geq2$ 的比例 & 43.75\% \\
    负向且 $z_{H1}\leq-2$ 的比例 & 39.06\% \\
    Size 月均系数 & $-0.4540\%$ \\
    Size 月均系数的 IID $t$ 值 & $-2.1758$ \\
    平均 RLM 降权比例 & 23.56\% \\
    平均伪 $R^2$ & 12.22\% \\
    \bottomrule
  \end{tabular}
\end{table}

EP 已按形成日横截面标准化，因此月均系数 0.1623\% 表示 EP 提高一个横截面
标准差时，下一持有期收益平均提高约 0.1623\%。简单乘以 12 后约为 1.95\%，
但这只是因子溢价尺度，不是可交易组合年化收益。

\subsection{分年度结果}

\begin{table}[H]
  \centering
  \caption{EP 月度横截面 RLM 分年度汇总}
  \begin{tabular}{crrrrr}
    \toprule
    年份 & 月数 & EP 月均系数 & IID $t$ 值 & IID $p$ 值 & 正系数比例 \\
    \midrule
    2020 & 6  & 0.2476\%  & 1.2853  & 0.2550 & 66.67\% \\
    2021 & 12 & $-0.0067\%$ & $-0.0274$ & 0.9787 & 58.33\% \\
    2022 & 12 & $-0.0323\%$ & $-0.1801$ & 0.8603 & 41.67\% \\
    2023 & 12 & 0.3872\%  & 1.7288  & 0.1118 & 75.00\% \\
    2024 & 12 & 0.4144\%  & 1.0109  & 0.3338 & 50.00\% \\
    2025 & 10 & $-0.0249\%$ & $-0.0970$ & 0.9248 & 30.00\% \\
    \bottomrule
  \end{tabular}
\end{table}

\subsection{结果解释}

EP 的全期平均系数方向符合预期，但在月度系数独立同方差假设下
$p=0.1524$，未达到常用的 5\% 显著性标准。单月 H1 $z$ 值显著的比例较高，
而全期均值不显著，原因是显著月份中正向占 43.75\%、负向占 39.06\%，方向
抵消明显。年度结果也不稳定：2023--2024 年为正，2021--2022 年及 2025 年
接近零或为负。因此，当前 RLM 结果只能说明 EP 存在正向平均溢价迹象，尚不能
单独认定其具有稳定选股能力，还需要结合 Rank IC、行业内分层回测以及 OLS/WLS
对照。

\begin{figure}[H]
  \centering
  \includegraphics[width=0.96\textwidth]{rlm_ep_premium.png}
  \caption{RLM 月度 EP 系数与累计因子溢价。累计因子溢价是月度系数的累计和，
  不是可交易组合净值。}
  \label{fig:rlm-premium}
\end{figure}

完整逐月结果、样本审计和回归权重保存在
\texttt{RLM\_H1\_独立同方差} 结果目录下的 \texttt{RLM统计结果.md} 文件中。

\end{document}
